\section{Responsabilidades, Limitaciones y Restricciones}

\subsection{Responsabilidades del Proveedor del Servicio}
\begin{itemize}
    \item Mantener la infraestructura de red en óptimas condiciones.
    \item Monitorear la red para detectar y resolver problemas de manera proactiva.
    \item Proporcionar soporte técnico de acuerdo a lo establecido en este SLA.
    \item Notificar a los usuarios sobre mantenimientos programados que puedan afectar el servicio.
\end{itemize}

\subsection{Responsabilidades de los Usuarios}
\begin{itemize}
    \item Utilizar la red de manera responsable y de acuerdo a las políticas de la universidad.
    \item No realizar actividades que puedan comprometer la seguridad o el rendimiento de la red.
    \item Reportar cualquier problema de red a través de los canales de comunicación establecidos.
    \item Proporcionar información clara y precisa al solicitar soporte técnico.
\end{itemize}

\subsection{Limitaciones y Restricciones}
\begin{itemize}
    \item El proveedor del servicio no se hace responsable por problemas causados por el mal uso de la red por parte de los usuarios.
    \item La disponibilidad del servicio puede verse afectada por factores fuera del control del proveedor, como cortes de energía eléctrica o fallas en el proveedor de servicios de Internet.
    \item Se podrán aplicar restricciones de ancho de banda a ciertos servicios o aplicaciones para garantizar un rendimiento equitativo para todos los usuarios.
\end{itemize}
